\documentclass[11pt]{article}
\usepackage{amsmath,amssymb,amsthm}
\newtheorem{theorem}{Theorem}
\newtheorem{proposition}{Proposition}
\newtheorem{corollary}{Corollary}
\begin{document}

\begin{theorem}[Exact finite-state anomaly criterion]
For any arrangement pattern $z$ with $P(Z=z)>0$ in the hierarchical finite-$M$
MSRC model, define
\[
J[z,i]=c_i(z)+(\pi_{\rm old}\otimes h_i)^T e^{A_4T_A}F_z.
\]
Then
\[
\Gamma(z)>0
\iff
\max\{J[z,2],J[z,3]\}>J[z,1],
\]
and
\[
\Delta(z)=\frac{J[z,2]-J[z,3]}{P(Z=z)}.
\]
\end{theorem}

\begin{proof}
The common positive denominator is
$P(Z=z)=\sum_iJ[z,i]$. Therefore comparisons among conditional probabilities
are exactly comparisons among their joint numerators. Subtracting the two
operator expressions gives the stated differences.
\end{proof}

\begin{theorem}[A zero-switching crossed-pattern bound]
Assume a fixed frequency $p\in(0,1)$, zero switching, chromosome count $M$,
and the finite-floor coalescence rates
\[
c_0=1/[M(1-p)],\qquad c_1=1/(Mp).
\]
For $Z=0101$, if $T_I>M(1-p)\log2$, then $Q_2$ is more probable than $Q_1$.
For $Z=0110$, if $T_I>Mp\log2$, then $Q_3$ is more probable than $Q_1$.
The complementary cases $1010$ and $1001$ obey the same bounds with the
matching background exchanged.
\end{theorem}

\begin{proof}
For $0101$, lineages 1 and 2 have different backgrounds and cannot coalesce on
the sister branch. On the intermediate branch, 1 and 3 are the only matching
pair among the three lineages, with rate $c_0$. Its coalescence probability is
$1-e^{-c_0T_I}$ and fixes $Q_2$. Hence
$P(Q_2\mid Z=0101)\ge1-e^{-c_0T_I}$, while $Q_1$ is impossible on this event
and so $P(Q_1\mid Z=0101)\le e^{-c_0T_I}$. The displayed condition makes the
first bound larger than the second. The other cases follow by exchanging the
matching pair and/or backgrounds.
\end{proof}

\begin{proposition}[Complement transformation]
Under neutral dynamics, replacing every arrangement bit by its complement and
replacing $p$ by $1-p$ preserves the conditional quartet probabilities after
mapping $z$ to its complement. With genic selection, the same statement holds
under the additional parameter transformation $s\mapsto-s/(1+s)$.
\end{proposition}

\begin{proof}
The background switching rates and finite-floor coalescence rates are
interchanged by the transformation. The Moran count generator is unchanged in
the neutral case. Under genic selection, exchanging the two relative fitnesses
and rescaling by the former A1 fitness gives the stated transformed selection
coefficient. The operator recursion is therefore conjugate under the bit
complementation map.
\end{proof}

\begin{proposition}[Fast-switching fixed-frequency limit]
At fixed $p$ and fixed finite branch durations, the leading
$\rho\to\infty$ averaged pair-coalescence rate is $1/M$. Thus arrangement
conditioning loses its leading-order effect and the ordinary MSC quartet law is
recovered at the averaged level.
\end{proposition}

\begin{proof}
The one-lineage switching chain has stationary probabilities $(1-p,p)$.
A pair is in the same background with probabilities $(1-p)^2$ and $p^2$.
Averaging the two hazards gives
$(1-p)^2/[M(1-p)]+p^2/(Mp)=1/M$. Standard finite-state fast-slow averaging
then gives the leading limit. A full coupled finite-M singular-perturbation
bound is not asserted here.
\end{proof}

The exact finite-M operator criterion is necessary and sufficient. The crossed
pattern bound is sufficient but not necessary; numerical roots in the full
operator are not claimed to be unique without an additional monotonicity
proof.
\end{document}
